\chapter{Conclusions and Future Work}
\label{ch:conclusion}

\section{Summary}

This thesis set out to make the SIMD vector accelerator of the Klessydra RISC-V
architecture both \emph{fault-tolerant} and \emph{observable in the field}, and
to do so in a way that is reusable and automatable rather than bespoke. Three
contributions delivered on that goal.

First, the originally monolithic accelerator was \emph{re-engineered into a clean
modular hierarchy} in which each functional unit --- adder, multiplier, shifter,
comparator, and accumulator --- is an independently instantiable VHDL entity with
an explicit interface. This restructuring is the enabling substrate for
everything that follows: it is what allows redundancy and instrumentation to be
attached at arbitrary granularity without disturbing the surrounding logic.

Second, a \emph{reusable fault-tolerance framework} was built on top of this
structure from two tiny, technology-independent primitives --- a DMR comparison
cell that provides error detection through an error flag, and a TMR majority
voter that provides error masking. Wrapping the functional units with these
primitives produced a complete family of hardened configurations: each unit
hardened individually, the full datapath hardened together, and an alternative
replication-based scheme that votes already-existing parallel lanes.

Third, an \emph{automated IEEE~1687 (IJTAG) instrumentation flow} was developed
that inserts on-chip sensors into the VHDL hierarchy from pragma markers and a
single JSON description, and generates the Tessent scripts that extract the
instrument connectivity and build the hierarchical access network with Segment
Insertion Bits --- giving in-field observability with no manual rewiring.

Finally, every configuration was synthesised across three technology nodes
spanning a mature, a mainstream, and an advanced process, and the resulting data
were reduced to a consistent set of comparisons. The exploration produced several
clear and actionable findings:

\begin{itemize}
  \item The cost of redundancy is dominated by \emph{which} functional unit is
        hardened: the multiplier is by far the most expensive and the comparator
        almost free, so fine-grained, selective hardening is highly effective.
  \item Hardening only the datapath keeps full-accelerator DMR and TMR well below
        the naive $2\times$ and $3\times$ factors, because the control, mapping,
        and memory-interface logic remain single-copy.
  \item Redundancy adds area and power but only a few percent of critical-path
        delay (at most about $1.06\times$, at the mature node), because the
        comparators and voters are shallow layers that rarely sit on the
        critical path.
  \item When the accelerator is already replicated for throughput, adding voting
        on top is nearly free in area --- redundancy economics are
        context-dependent.
  \item The in-field monitoring instrumentation costs at most a couple of percent
        and is independent of the redundancy choice.
  \item All of these overhead factors are \emph{stable across the technology
        range}, so the resulting design guidelines transfer between nodes.
\end{itemize}

Together, these results constitute a reusable, automatable methodology for
hardening a RISC-V vector accelerator against faults while keeping it observable
in the field, accompanied by quantitative guidelines for navigating the
protection-versus-cost trade-off.

\section{Limitations}

Several limitations should be kept in mind when interpreting the results. The
power figures are estimates from default switching activity rather than from
gate-level activity captured on representative workloads, so they are most
meaningful as ratios. Timing was characterised against a single aggressive clock
target that no configuration meets, so the comparison rests on relative
critical-path (data-arrival) delays rather than on closing timing at a realistic
frequency. The fault-tolerance claims are \emph{structural}:
DMR detects and TMR masks any single discrepancy by construction, but the study
did not run an explicit fault-injection campaign to measure the resulting
reduction in failure rate. Finally, the advanced-node replication runs collapsed
to an identical netlist and were therefore not over-interpreted.

\section{Future work}

The framework opens several natural directions.

\begin{itemize}
  \item \textbf{Fault-injection validation.} Coupling the hardened
        configurations with an RTL or gate-level fault-injection campaign would
        turn the structural guarantees into measured reliability metrics (e.g.,
        architectural vulnerability factor and failures-in-time), and would let
        selective-hardening decisions be guided by measured per-unit
        vulnerability rather than area alone.
  \item \textbf{Constrained and gate-level-activity synthesis.} Repeating the
        sweep with binding clock constraints and workload-derived switching
        activity would refine the power and timing figures and reveal any
        frequency penalty at aggressive targets.
  \item \textbf{Richer redundancy.} Triplicated voters for the most critical
        outputs, residue or parity checking for the multiplier, and selective
        ECC on the scratchpad would broaden the protection palette and could be
        evaluated within the same modular structure.
  \item \textbf{Closed-loop monitoring.} The IJTAG activity and thermal sensors
        could drive run-time adaptation --- throttling, voltage/frequency
        scaling, or reconfiguration of the redundancy mode --- turning passive
        observability into active in-field health management. Adding further
        sensor types (thermal, ageing, voltage) to the existing tool is
        straightforward given its sensor-agnostic interface.
  \item \textbf{Broader design space.} Sweeping the SIMD width and lane count
        together with the redundancy options would map how the cost structure
        scales with accelerator size, and extending the flow to other Klessydra
        cores or to other RISC-V accelerators would test the generality of the
        methodology.
\end{itemize}

In bringing structural redundancy and standardised in-field monitoring together
on one parameterised RISC-V vector accelerator, and characterising the result
across three technology nodes, this thesis provides both a practical toolset and
an empirical foundation for designing dependable RISC-V acceleration in
safety- and mission-critical systems.
